% TOP_PROBLEM: {"id": 3233, "title": "Hedetniemi's Conjecture", "category_id": 3, "category": {"id": 3, "name": "graph_theory", "display_name": "Graph Theory", "description": "Problems involving graphs, networks, and their properties.", "slug": "graph-theory", "order_index": 3, "created_at": "2026-07-31T15:26:25.671Z"}, "status": "open"}
% FIRST_POSTED: null
% ENTRY_KIND: statement_only
% Imported from: batch02.tex; UnsolvedMath ID 3233
% Compile twice: lualatex 3233.tex
\documentclass[11pt]{article}
\usepackage{fontspec}
\setmainfont{Latin Modern Roman}
\usepackage{amsmath,amssymb,mathtools,unicode-math}
\setmathfont{Latin Modern Math}
\usepackage{xurl,hyperref}
\hypersetup{hidelinks,pdftitle={UnsolvedMath Problem 3233}}
\newcommand{\sref}[2]{\hyperlink{source:#1:#2}{[S#2]}}
\newcommand{\cataloguescope}[1]{\paragraph{Mathematical question.}#1}
\begin{document}
% UnsolvedMath ID 3233; source catalogue code OPG-806
\setcounter{section}{3232}
\section{Hedetniemi's Conjecture}\label{3233}
{\small\sffamily UnsolvedMath ID 3233 · Graph Theory}
\subsection{Definitions and mathematical statement}

\paragraph{Shared notation.}$\mathbb N=\{1,2,\ldots\}$, and $\mathbb Z,\mathbb Q,\mathbb R,\mathbb C$ denote the integers, rationals, reals and complex numbers. Any other conventions are specified locally.


For finite simple graphs $G,H$, the tensor product $G\times H$ has vertices $(g,h)$ and joins $(g,h)$ to $(g',h')$ exactly when $gg'\in E(G)$ and $hh'\in E(H)$. Write $G\to H$ for a graph homomorphism. For $n\ge2k>0$, let $K_{n/k}$ have vertices $\mathbb Z/n\mathbb Z$, with distinct vertices adjacent when their cyclic distance is at least $k$; define $\chi_c(G)=\inf\{n/k:G\to K_{n/k}\}$ for graphs with an edge, and $\chi_c(G)=1$ for nonempty edgeless graphs. The exponential graph $K_n^G$ has as vertices all maps $V(G)\to V(K_n)$, with $f,f'$ adjacent when $f(x)f'(y)\in E(K_n)$ for every ordered pair $(x,y)$ with $xy\in E(G)$. A graph $K$ is multiplicative if $G\times H\to K$ always implies $G\to K$ or $H\to K$.
\paragraph{Statement recorded in the supplied source.}
The recorded conjecture is \[\chi(G\times H)=\min\{\chi(G),\chi(H)\}.\] Equivalently, whenever $\chi(G)>n$, must $K_n^G$ be $n$-colorable? The source also records Zhu's stronger circular version \[\chi_c(G\times H)=\min\{\chi_c(G),\chi_c(H)\},\] equivalently multiplicativity of every circular complete graph $K_{n/k}$. It additionally defines $f(n)=\min\{\chi(G\times H):\chi(G)=\chi(H)=n\}$, so the original equality is $f(n)=n$. \sref{3233}{2}
\subsection{Short English statement}
Does the tensor product preserve the smaller chromatic number? The source also asks the circular-coloring version and its equivalent homomorphism formulations.
\subsection{Sources}
\begin{description}
\item[\hypertarget{source:3233:1}{S1}] ULAM.AI and the UnsolvedMath Contributors, UnsolvedMath: A Curated Collection of Open Mathematics Problems, record 3233 (OPG-806). CC BY 4.0. \url{https://huggingface.co/datasets/ulamai/UnsolvedMath}.
\item[\hypertarget{source:3233:2}{S2}] Open Problem Garden, Hedetniemi's Conjecture, original node 806, statement and mathematical discussion. \url{https://www.openproblemgarden.org/op/hedetniemis_conjecture}.
\end{description}
\label{end:3233}
\end{document}
